% ======================================================
%  Rozdział 1 — Arytmetyka: Podstawy
%
%  SAMPLE. Three exercises, ported by hand from areas/1-arytmetyka-podstawy.tex
%  to prove the frame machinery on real content. The remaining seven, and the
%  other four areas, are converted in issue #17 -- script-assisted, because the
%  deck macros are regular enough to convert mechanically and re-run while the
%  frame design settles.
%
%  THE PATTERN. One exercise per frame. The answer to an exercise opens the
%  frame AFTER it, so the reader covers the page below the hairline, works the
%  exercise against the stopwatch, then reads on to check. Every exercise
%  therefore appears exactly once and every answer exactly once, and the last
%  frame of a chapter carries only the final answer.
%
%  Diacritics are written as UTF-8 here, not as the deck's \. z escapes. Both
%  reach the same glyph under inputenc; the conversion script must preserve
%  whichever form it finds rather than normalising one into the other.
% ======================================================

\chapter{Arytmetyka: Podstawy}

Dodawanie, odejmowanie, mnożenie i dzielenie — na czas i bez kalkulatora.
Zakryj dłonią stronę poniżej kreski, uruchom stoper, rozwiąż zadanie, a
dopiero potem czytaj dalej.

% -- 1 --
\begin{fr}
  \exercise[\CategoryBadge[ArithColor!20]{Dodawanie}]{1}{30 s}{$47 + 38 = ?$}
\end{fr}

% -- 2 --
\begin{fr}
  \ans{85}{}
  \exercise[\CategoryBadge[ArithColor!20]{Odejmowanie}]{1}{30 s}{$156 - 79 = ?$}
\end{fr}

% -- 3 --
\begin{fr}
  \ans{77}{}
  \exercise[\CategoryBadge[ArithColor!20]{Mnożenie}]{1}{45 s}{$24 \times 7 = ?$}
\end{fr}

% The chapter's last frame carries the last answer and nothing else. A hint is
% shown here to exercise the second argument of \ans on a real page; most
% exercises in this area need none.
\begin{fr}
  \ans{168}{$24 \times 7 = 20 \times 7 + 4 \times 7$}
\end{fr}
